% year: תשפ"ב
% type: מבחן
% semester: א
% moed: א
% number: 3
% points: 17
% categories: taylor
% source: exams/מבחנים/תשפב/מועד א תשפב.pdf

\begin{question}
(17 נק') העריכו בעזרת קירוב לינארי (פולינום טיילור ממעלה ראשונה) את $\sqrt[4]{18}$ ותנו חסם לשגיאה.
\end{question}

\begin{hint}
פתחו את $f(x)=\sqrt[4]{x}$ סביב נקודה קרובה ל-$18$ שבה השורש הרביעי ידוע.
\end{hint}

\begin{hint}
את השגיאה חסמו בעזרת שארית לגרנז': $R_1(x)=\frac{f''(c)}{2}(x-a)^2$ עבור $c$ בין $a$ ל-$x$.
\end{hint}

\begin{solution}
נבחר $f(x)=x^{1/4}$ ונפתח סביב $a=16$ (כי $\sqrt[4]{16}=2$).
\[ f'(x)=\tfrac14x^{-3/4},\qquad f''(x)=-\tfrac{3}{16}x^{-7/4}. \]
לכן $f(16)=2$, $f'(16)=\frac14\cdot16^{-3/4}=\frac14\cdot\frac18=\frac1{32}$, ופולינום טיילור מסדר ראשון הוא
\[ P_1(x)=2+\frac{1}{32}(x-16). \]
הקירוב:
\[ \sqrt[4]{18}\approx P_1(18)=2+\frac{2}{32}=\frac{33}{16}=2.0625. \]

\textbf{חסם לשגיאה.} לפי משפט טיילור עם שארית לגרנז' ($f$ גזירה פעמיים ב-$(0,\infty)$), קיימת $c\in(16,18)$ כך ש-
\[ R_1(18)=f(18)-P_1(18)=\frac{f''(c)}{2!}(18-16)^2=-\frac{3}{16}c^{-7/4}\cdot\frac{4}{2}=-\frac{3}{8}c^{-7/4}. \]
הפונקציה $c^{-7/4}$ יורדת, ולכן עבור $c>16$:
\[ |R_1(18)|=\frac38c^{-7/4}<\frac38\cdot16^{-7/4}=\frac38\cdot\frac{1}{2^7}=\frac{3}{1024}\approx0.0029. \]

\textbf{תשובה:} $\sqrt[4]{18}\approx2.0625$ עם שגיאה קטנה מ-$\frac{3}{1024}\approx 0.003$. יתרה מזו, כיוון ש-$R_1<0$, הקירוב גדול מהערך האמיתי:
$2.0625-\frac{3}{1024}<\sqrt[4]{18}<2.0625$.
(לבדיקה: $\sqrt[4]{18}\approx2.05977$.)
\end{solution}
